\documentclass[10pt,a4paper]{article}
\usepackage[utf8]{inputenc}
\usepackage[T1]{fontenc}
\usepackage[margin=1.5cm]{geometry}
\usepackage{hyperref}
\usepackage{enumitem}
\usepackage{xcolor}
\usepackage{titlesec}

% Palette de couleurs sobre et professionnelle
\definecolor{primary}{RGB}{30, 58, 138}    % Bleu marine élégant
\definecolor{textdark}{RGB}{31, 41, 55}   % Gris très foncé (plus doux que noir pur)
\definecolor{linegray}{RGB}{209, 213, 219}

\color{textdark}
\pagestyle{empty}

% Configuration des liens hypertexte
\hypersetup{
    colorlinks=true,
    linkcolor=primary,
    urlcolor=primary,
    pdftitle={CV - Badroudine Bacar},
    pdfauthor={Badroudine Bacar}
}

% Configuration des titres de sections (compatible ATS)
\titleformat{\section}
  {\large\bfseries\color{primary}\uppercase}
  {}{0em}{}[\vspace{-0.5em}\rule{\textwidth}{0.8pt}\vspace{-0.3em}]

\titlespacing*{\section}{0pt}{1.0em}{0.6em}

\setlist[itemize]{leftmargin=1.2em, itemsep=2pt, topsep=2pt, parsep=0pt}

\begin{document}

% --- EN-TÊTE ---
\begin{center}
    {\huge \textbf{\color{primary} Badroudine Bacar}} \\[0.4em]
    {\large \textbf{AI / Machine Learning Engineer (Stage PFE)}} \\[0.4em]
    \small
    Calais / Hauts-de-France (Mobilité : Suisse, Luxembourg, Belgique, Télétravail) \ $|$ \ 
    \href{mailto:badroudinebacar@gmail.com}{badroudinebacar@gmail.com} \ $|$ \ 
    \href{https://www.linkedin.com/in/badroudine-bacar/}{LinkedIn} \ $|$ \ 
    \href{https://github.com/Badrouuu}{GitHub}
\end{center}

\vspace{-0.4em}
\noindent
\small Élève-ingénieur en dernière année à l'EILCO (spécialité IA), doté d'une certification TOEIC 905/990. Spécialisé en Machine Learning, NLP et Deep Learning, avec une expérience avérée dans le développement de pipelines RAG, d'agents LLM (Few-Shot Prompting) et de calculs d'espaces vectoriels sémantiques (Word2Vec). Expert dans la conception d'APIs haute performance avec FastAPI et leur conteneurisation via Docker. À la recherche d'un stage de fin d'études (PFE) de 6 mois à partir de mars 2027.

% --- FORMATION ---
\section{Formation \& Certifications}
\begin{itemize}

    \item \textbf{Diplôme d'Ingénieur Informatique — Parcours Intelligence Artificielle} \hfill 2024 - 2027 \\
    \textit{EILCO (École d'Ingénieurs du Littoral Côte d'Opale), Calais, France} \\
    \footnotesize Dernière année (ING3). Modules avancés : Deep Learning, Vision par Ordinateur, Machine Learning, Architecture Logicielle, MLOps, CI/CD, Bases de Données Avancées.

    \item \textbf{CPGE TSI (Classe Préparatoire aux Grandes Écoles)} \hfill 2022 - 2024 \\
    \textit{Lycée Antonin Artaud, Marseille, France} \\
    \footnotesize Formation intensive pluridisciplinaire : Mathématiques fondamentales, Physique appliquée, Sciences de l'Ingénieur et Algorithmique.

    \item \textbf{Certification Anglais : TOEIC Listening and Reading} \hfill Score : \textbf{905 / 990 (Niveau C1)}
\end{itemize}

% --- COMPÉTENCES TECHNIQUES ---
\section{Compétences Techniques}
\begin{itemize}
    \item \textbf{Intelligence Artificielle \& Data :} PyTorch, Scikit-Learn, Gensim (Word Embeddings), Google ML Kit (OCR / Vision), OpenCV, RAG \& LLM Prompting, NumPy, Pandas
    \item \textbf{Langages de Programmation :} Python, TypeScript, JavaScript, SQL, C, Java, Dart, Bash
    \item \textbf{Backend, Cloud \& DevOps :} FastAPI, Docker, Linux, Git \& GitHub, Supabase (PostgreSQL), Firebase Cloud Functions, CI/CD basics, RESTful APIs
    \item \textbf{Langues :} Français (Langue maternelle) $|$ Anglais (Niveau C1 certifié — TOEIC 905/990)
\end{itemize}

% --- EXPÉRIENCE PROFESSIONNELLE ---
\section{Expérience Professionnelle}

\noindent
\textbf{Fondateur \& Développeur Lead (Stage Entrepreneurial PÉPITE / CEL)} \hfill 2025 - 2026 \\
\textit{Étiquette Jaune (Micro-entreprise) \ --- \ Calais, France}
\begin{itemize}

    \item Accompli l'intégration de modèles de Vision par Ordinateur mesurés par la vitesse d'extraction textuelle, en faisant usage de Google ML Kit Text Recognition (OCR).

    \item Accompli le développement d'un backend serverless mesuré par la synchronisation temps réel, en faisant appel à Google Cloud Firestore, Cloud Functions et TypeScript.

    \item Accompli la mise en place de modules critiques mesurés par la sécurité des transactions, en faisant appel à Firebase Auth et l'API Stripe.

    \item \textbf{Technologies :} Flutter, Dart, Google ML Kit (OCR), Firebase, Cloud Functions (TypeScript), Stripe API, Git
\end{itemize}
\vspace{0.4em}


% --- PROJETS CLÉS ---
\section{Projets Techniques \& Réalisations}

\noindent
\textbf{VectorWord / Guess-What (NLP \& Vector Space Search)} \hfill 2026 \\
\textit{Développeur Fullstack \& ML}
\begin{itemize}

    \item Accompli la mise en place d'une API haute performance mesurée par un temps de réponse optimisé, en faisant appel à FastAPI et Uvicorn pour servir des calculs de similarité cosinus sur des espaces vectoriels Word2Vec multilingues (FR/EN) via Gensim et NumPy.

    \item Accompli la portabilité et le déploiement industrialisé mesurés par une reproductibilité totale, en faisant usage d'une architecture conteneurisée avec Docker et un déploiement automatisé sur Railway.

    \item Accompli le développement d'un frontend réactif mesuré par une synchronisation temps réel des données, en faisant appel à Next.js 15, TypeScript, Tailwind CSS et Supabase (PostgreSQL).

    \item \textbf{Technologies :} Python, FastAPI, Gensim (Word2Vec), NumPy, Docker, Supabase, Next.js 15, TypeScript
\end{itemize}
\vspace{0.4em}

\noindent
\textbf{Pipeline RAG \& Automatisation de Code (Agentic AI)} \hfill 2026 \\
\textit{Ingénieur RAG / LLM}
\begin{itemize}

    \item Accompli la conception d'un moteur de parsing d'offres et de matching mesuré par la précision du score de pertinence sémantique, en faisant usage d'un système agentique RAG et LLM.

    \item Accompli la génération automatisée de documents mesurée par la réduction drastique des hallucinations, en faisant appel à des stratégies de Few-Shot Prompting et de génération structurée LaTeX.

    \item \textbf{Technologies :} Python, LLMs, RAG, Prompt Engineering, LaTeX, NLP
\end{itemize}
\vspace{0.4em}

\noindent
\textbf{Computer Vision \& Deep Learning Lab (EILCO)} \hfill 2026 \\
\textit{Projets académiques ING3}
\begin{itemize}

    \item Accompli l'entraînement et l'optimisation de modèles mesurés par la performance de classification et segmentation, en faisant appel à PyTorch.

    \item Accompli le prétraitement et l'extraction de caractéristiques mesurés par la robustesse des données, en faisant usage d'OpenCV, Scikit-Learn et NumPy.

    \item \textbf{Technologies :} Python, PyTorch, OpenCV, Scikit-Learn, Matplotlib
\end{itemize}
\vspace{0.4em}


\end{document}